\section{Perspectivas 2026: direcciones futuras del aprendizaje por refuerzo}

\subsection{Tareas más difíciles, entornos más abiertos}

Las perspectivas de Cui Ganqu (崔干曲) giran en torno a la palabra clave «Goal-Oriented»: la característica esencial del aprendizaje por refuerzo es que está orientado a objetivos, a diferencia del paradigma de clonación de conducta del Pretraining y el SFT. Cuando el modelo necesita realizar tareas para las que recolectar datos es extremadamente difícil, es necesario cambiar a este método orientado a objetivos que es el RL.

\begin{importantbox}{Las líneas principales del RL para 2026 (Cui Ganqu)}
\begin{enumerate}
    \item \textbf{Tareas más difíciles}: el modelo necesita interactuar durante periodos largos dentro del Sandbox, prestando atención solo al resultado final
    \item \textbf{Escenarios más abiertos}: entornos externos más complejos, fuentes de Feedback más variadas
    \item \textbf{Interacciones de mayor alcance}: sistemas de Memory, Self-Evolve (autoevolución)
    \item \textbf{Los retos de Infra y arquitectura}: este aumento de los grados de libertad traerá enormes retos tanto a nivel de Infra como de arquitectura del modelo
\end{enumerate}
\end{importantbox}

\subsection{Scaling: el doble cuello de botella de Infra y datos}

Hu Jian (胡建) considera que el núcleo del RL en 2026 sigue siendo el problema de \textbf{Scaling}:

\begin{enumerate}
    \item \textbf{Infra Scaling}: cómo lograr un entrenamiento estable, a gran escala y de largo plazo
    \item \textbf{Scaling de datos}: el Agentic RL necesita una gran cantidad de entornos de alta calidad (reales o sintéticos); cómo ampliar la escala de los entornos es el cuello de botella clave: «si siempre se aprende con pocos datos, el límite superior está ahí»
    \item \textbf{Convergencia de los algoritmos}: es posible que ya no haya grandes cambios a nivel de los algoritmos de RL, sino más bien innovaciones puntuales (como el Process-based Value Model)
\end{enumerate}

\subsection{RL totalmente asíncrono: el siguiente campo de batalla de la industria}

Zheng Chujie (郑楚杰) hizo un pronóstico técnico concreto:

\begin{importantbox}{La dirección del RL industrial en 2026 (Zheng Chujie)}
El \textbf{RL totalmente asíncrono (Fully Asynchronous RL)} se convertirá en la corriente principal. La contradicción central es el equilibrio entre \textbf{la utilización del cómputo y el desempeño del algoritmo}:
\begin{itemize}
    \item En un marco totalmente asíncrono, una misma respuesta del modelo puede ser generada por varias versiones del modelo, lo que introduce de forma natural el problema de Off-Policy
    \item Es necesario estudiar la adaptación de los algoritmos a escenarios asíncronos: dada una receta de algoritmos, ¿cuál es el grado de tolerancia al Off-Policy?
    \item El soporte actual de los marcos de código abierto para el RL totalmente asíncrono todavía no es maduro; este será el trabajo prioritario de los próximos meses
    \item Primero hay que elevar la eficiencia de cómputo, y solo entonces se podrá discutir el problema del Scale
\end{itemize}
\end{importantbox}

\subsection{Agentic RL y el Credit Assignment de largo alcance}

Las perspectivas de Li Yingru (李英儒) se centran en el \textbf{Agentic RL} y en los \textbf{problemas de largo alcance}:

\begin{enumerate}
    \item \textbf{La interacción con el entorno se convierte en el nuevo Bottleneck}: al pasar del Reasoning RL al Agentic RL, el cuello de botella se desplaza de la GPU Generation a la interacción con el entorno (esperar la compilación, la retroalimentación del Sandbox, etcétera); el problema de Sample Efficiency se vuelve más notorio
    \item \textbf{El Effective Horizon crece de forma drástica}: un Agent puede necesitar «programar durante tres días» para completar una tarea; los verdaderos puntos de decisión de bifurcación (el 20\% de los nodos clave) dan lugar a un Effective Horizon extremadamente largo
    \item \textbf{El Credit Assignment se convierte en un reto clave}: en escenarios Agentic de largo alcance, estimar la Token-Level Advantage a partir únicamente de la Outcome Reward puede dejar de ser suficiente; se necesitan métodos de Credit Assignment y Variance Reduction de grano más fino
\end{enumerate}

\begin{knowledgebox}{Aprender de la sabiduría del RL tradicional}
Li Yingru insiste repetidamente en que muchos de los problemas que enfrenta el RL de los modelos grandes ya han sido estudiados a fondo en el RL tradicional. Los requisitos de marco del Agentic RL se parecen cada vez más a los del RL de videojuegos (interacción entre el Agent y el entorno, muestreo asíncrono, arquitectura desacoplada), y temas clásicos como el Sample Efficiency, el Credit Assignment y el Variance Reduction cobran nueva vida en estos nuevos escenarios.
\end{knowledgebox}

\subsection{Resumen de la sección}

El desarrollo del RL en 2026 avanzará a lo largo de tres líneas principales: (1) la complejidad de las tareas seguirá aumentando, pasando del Reasoning al Agentic; (2) el Scaling de Infra y de datos se convertirá en el cuello de botella central, y el RL totalmente asíncrono será la corriente principal en la industria; (3) problemas clásicos del RL como el Credit Assignment de largo alcance recuperarán atención en los nuevos escenarios. Es posible que las grandes transformaciones a nivel algorítmico se desaceleren, pero todavía queda un amplio espacio para la optimización continua del Recipe y para su explicación teórica.

